% COSC 483 Cryptographic Playground report -- ACM Primary Article Template (acmart, sigconf)
% Build: latexmk -pdf main.tex   (or upload this folder to Overleaf, which ships acmart)
\documentclass[sigconf,nonacm]{acmart}

\title{A Signal-Style End-to-End Encrypted Messenger: X3DH and the Double Ratchet over an Untrusted Relay}

\author{Liam Loeffel}
\affiliation{\institution{University of Tennessee, Knoxville}\country{USA}}
% \author{Teammate}\affiliation{\institution{University of Tennessee, Knoxville}\country{USA}}

\begin{abstract}
TODO (write last).
\end{abstract}

\begin{document}
\maketitle

% ---------------------------------------------------------------- 20 pts
\section{Introduction}
% General area -> specific problem -> our approach -> what's different from others.

% ---------------------------------------------------------------- 20 pts
\section{Background}
% Threat model of a relay server; forward secrecy vs. post-compromise security;
% X3DH~\cite{x3dh}; Double Ratchet~\cite{doubleratchet}; primitives (X25519~\cite{rfc7748},
% Ed25519~\cite{rfc8032}, HKDF~\cite{rfc5869}, AES-GCM); formal analysis~\cite{cohn2017formal}.

% ---------------------------------------------------------------- 20 pts
\section{System Design}
% Architecture figure: client / relay server / prekey store. Message flow for first
% contact and steady state. Trust boundaries. Deviations from Signal (separate
% Ed25519 identity key instead of XEdDSA).

% ---------------------------------------------------------------- 20 pts
\section{Implementation and Testing}
% Module layout, key data structures, state persistence.
% Correctness: RFC known-answer tests, cross-check vs. a reference implementation,
% out-of-order/dropped/tampered/replayed messages.
% Security demos: forward secrecy, post-compromise security, server sees only ciphertext.
% Table of test results.

% ---------------------------------------------------------------- 20 pts
\section{Discussion}
% Lessons learned, limitations (no metadata privacy, trust on first use, single device),
% future work (safety numbers, groups/MLS, PQXDH).

\section*{Use of AI}
% Required: how AI was used, and how we checked its output. Transcripts are in transcripts.zip.

\bibliographystyle{ACM-Reference-Format}
\bibliography{references}
\end{document}
